For \(K=1\) and \(T=5\), write the three maximal minors as
\[
F_1=x_2^2-x_1x_3,\qquad
F_2=x_2x_3-x_1x_4,\qquad
F_3=x_3^2-x_2x_4.
\]
At \(x^\star=(1,-1,1,-1)\), their gradients are
\[
\nabla F_1=(-1,-2,-1,0),\quad
\nabla F_2=(1,1,-1,-1),\quad
\nabla F_3=(0,1,2,1),
\]
and \(\nabla F_2=-\nabla F_1-\nabla F_3\). The first and third vectors are independent, so the Jacobian rank is two. Also
\[
M(x^\star)=
\begin{pmatrix}-1&1\\1&-1\\-1&1\end{pmatrix},
\]
whose right kernel is spanned by \((1,1)'\). Thus both targets have nonzero alignment with \(v_0=2^{-1/2}(1,1)'\).

Choose \(\rho>0\) small enough that \((F_1,F_3)\) has rank two throughout \(U_\rho\), the local rank-one locus is its zero set, and both alignment numerators are positive there. The local change of variables supplied by the two independent gradients gives coordinates \((y,z)\in\mathbb R^2\times\mathbb R^2\) in which this locus is \(\{y=0\}\). Since \(Q=F_1^2+F_2^2+F_3^2\) and \(F_2\) is a smooth function of \((F_1,F_3,z)\),
\[
Q(y,z)=y'R(z)y+O(\|y\|_2^3),
\]
with \(R(z)\) uniformly positive definite after decreasing \(\rho\). At \(x^\star\), \(M'M\) has nonzero eigenvalue six and normalized null vector \(v_0\), whence
\[
A_k(x^\star)=e_k'\operatorname{adj}(M'M)e_k
=6(e_k'v_0)^2=3,
\qquad k=0,1.
\]
Continuity therefore bounds both \(A_k(y,z)\) above and away from zero on the smaller cube.

The pseudo-inverse identity gives
\[
\{H_k>u\}=\{Q/A_k<u^{-2}\}.
\]
For each tangential \(z\), the normal section is an ellipse with area
\[
u^{-2}\frac{\pi A_k(0,z)}{\sqrt{\det R(z)}}+o(u^{-2}).
\]
The remainder is uniform away from lower-dimensional intersections with the cube boundary, and the coordinate Jacobian is continuous and bounded above and away from zero. The boundary intersections have zero tangential measure. Integrating the uniformly bounded section areas therefore gives
\[
u^2P\{H_k(X)>u\mid X\in U_\rho\}
\longrightarrow C_{k,\rho},
\]
where \(C_{k,\rho}\) is the constant-density normalization times the integral of the displayed strictly positive continuous coefficient over the rank-one locus in the cube. Consequently \(0<C_{k,\rho}<\infty\) for both targets.

For numerical verification, a scrambled Sobol rule with seed \(1729\), \(2^{20}\) points, and \(\rho=0.05\), evaluating the diagonal entries of \((M'M)^{-1}\), gave
\[
\begin{array}{c|rrrrr}
u&40&60&80&100&200\\ \hline
u^2\widehat P(H_0>u)&174.36&178.10&179.74&180.26&181.12\\
u^2\widehat P(H_1>u)&174.36&178.06&179.73&180.10&181.66
\end{array}
\]
which verifies stabilization at a finite positive constant for both coordinates.

For \(K=1,T=4\),
\[
M=\begin{pmatrix}x_2&x_1\\x_3&x_2\end{pmatrix},
\qquad \det M=x_2^2-x_1x_3.
\]
Direct inversion of \(M'M\) gives
\[
H_0=\frac{\sqrt{x_1^2+x_2^2}}{|x_2^2-x_1x_3|},
\qquad
H_1=\frac{\sqrt{x_2^2+x_3^2}}{|x_2^2-x_1x_3|}.
\]
At a smooth target-active point the determinant is a transverse coordinate and the relevant numerator is bounded away from zero. Integrating first in that transverse coordinate gives a tail proportional to \(u^{-1}\), with no logarithm, proving \((\alpha_{\Sigma,k},m_{\Sigma,k})=(1,0)\).